\documentclass[a4paper]{jsarticle}
\usepackage{amsmath,amssymb}
\usepackage{graphicx}
\usepackage{xcolor}
\usepackage{hyperref}
\usepackage{enumitem}
\usepackage{listings}
\usepackage{booktabs}
\usepackage{amsthm}
\usepackage{siunitx}

\newtheorem{kimari}{決まり}[section]

\newcommand{\R}{\mathbb{R}}
\newcommand{\shelf}[1]{\textbf{棚#1}}
\DeclareMathOperator{\argmax}{arg\,max}
\definecolor{ai}{HTML}{2B4C6F}

\title{星見図書室の運営メモ}
\author{架空の司書 \and 架空の助手}
\date{2026年10月}

\begin{document}
\maketitle

\begin{abstract}
星見図書室の貸し出しと棚の並べ方を、印刷して配るためにまとめた架空の文書です。
\end{abstract}

\tableofcontents

\section{はじめに}\label{sec:intro}
この文書は \LaTeX{} で書かれています。図書室の棚は\shelf{A}から順に並びます%
（コメントは印刷されません）。
詳しくは\ref{sec:rules}節と式~\eqref{eq:books}を見てください\footnote{脚注は文書の最後にまとめて刷られます。}。

\section{貸し出しの決まり}\label{sec:rules}
\subsection{冊数}
一度に借りられる冊数 $n$ は、次の式で決めます。
\begin{equation}
  n = \min\left(5, \left\lfloor \frac{d}{7} \right\rfloor + 2\right) \label{eq:books}
\end{equation}
ここで $d$ は利用した日数です。実数の集合 $\R$ や $\argmax_{x \in \R} f(x)$ も書けます。

\begin{kimari}[返却]\label{rule:return}
借りた本は \qty{14}{\day} 以内に返す。書庫は \qty{20}{\degreeCelsius}、湿度 \qty{55}{\percent} に保つ。
\end{kimari}

\begin{proof}[理由]
次の人が待っているため。棚の総延長は \qty{1.2e3}{\metre} あり、探すのに時間がかかる\cite{hoshimi2026}。
\end{proof}

\begin{align*}
  a &= b + c \\
  d &= e + f
\end{align*}

\subsection{手順}
\begin{enumerate}
  \item 窓口で図書カードを見せる
  \item 本を選ぶ
  \begin{enumerate}[label=(\alph*)]
    \item 雑誌は二冊まで\label{item:magazine}
    \item 辞書は館内だけ
  \end{enumerate}
  \item 返す日を確かめる
\end{enumerate}
雑誌の決まりは(\ref{item:magazine})です。

\begin{itemize}
  \item \textcolor{ai}{藍色の文字}と\emph{強調}と\texttt{等幅}
  \item[\checkmark] チェックの印を付けた項目
\end{itemize}

\begin{description}
  \item[貸し出し] 一度に五冊まで
  \item[予約] 窓口で受け付けます
\end{description}

\section{棚の一覧}
\begin{table}[h]
  \centering
  \caption{棚と分類}\label{tab:shelves}
  \begin{tabular}{lcr}
    \toprule
    棚 & 分類 & 冊数 \\
    \midrule
    A & 文学 & 1200 \\
    B & 科学 & 860 \\
    \multicolumn{2}{l}{合計} & 2060 \\
    \bottomrule
  \end{tabular}
\end{table}
表~\ref{tab:shelves}のとおりです。

\begin{lstlisting}[language=Python]
def lend(books, limit=5):
    return books[:limit]
\end{lstlisting}

\begin{tikzpicture}
  \draw (0,0) -- (1,1);
\end{tikzpicture}

\begin{quote}
静かにお過ごしください。---星見図書室
\end{quote}

詳しくは \url{https://example.com/hoshimi} か
\href{https://example.com/hoshimi/faq}{よくある質問}へ。決まり~\ref{rule:return}の背景は文献\cite{shelves,hoshimi2026}にあります。

\begin{thebibliography}{9}
\bibitem{hoshimi2026} 架空の司書『星見図書室の十年』星見出版, 2026.
\bibitem{shelves} 架空の助手「棚の並べ方の研究」『架空図書館学』第3巻, pp.~1--10, 2025.
\end{thebibliography}

\end{document}
